\documentclass[UTF8,a4paper,10pt]{ctexart}
\usepackage[margin=2.1cm]{geometry}
\usepackage{amsmath,amssymb,booktabs,hyperref}
\hypersetup{colorlinks=true,linkcolor=black,urlcolor=blue}
\title{面向 AnimalCLEF 2026 的开放集动物个体发现：\\全局表征融合与局部斑纹匹配}
\author{Fangwenky}
\date{2026 年 9 月}

\begin{document}
\maketitle

\begin{abstract}
AnimalCLEF 2026 要求将四种动物的测试照片按个体分簇；前三种动物有带身份标签的训练图像，德州角蜥则只有测试图像。本文比较 MegaDescriptor、MiewID、双模型特征融合、SIFT 局部匹配以及互近邻图聚类。以身份互斥的三种随机划分验证，双模型融合加蝾螈局部匹配使蝾螈保留集调整兰德指数（ARI）由基线的 $0.097\pm0.010$ 提高到 $0.222\pm0.023$；蠵龟由 $0.496\pm0.156$ 提高到 $0.835\pm0.050$。针对无标签角蜥，检查候选对的局部匹配稀疏性后采用更保守的聚类阈值。最终 Kaggle 提交的 public/private ARI 为 $0.30642/0.39072$，相较 MegaDescriptor 基线的 $0.13528/0.18564$ 明显提高。角蜥阈值经过排行榜比较，故该成绩不能作为独立泛化估计。
\end{abstract}

\section{问题定义}
给定带个体身份标签的参考图片集合 $R$ 与无标签测试集合 $Q$，目标是为每张测试图片 $x_i\in Q$ 赋予簇标签 $\hat y_i$。若两张图片展示同一个体，其簇标签应相同；不同个体应分属不同簇。对于前三种动物，测试个体可能在参考集合中出现，也可能完全未知；德州角蜥仅需在测试集合内部发现个体。比赛使用调整兰德指数（Adjusted Rand Index, ARI）比较预测划分与真实划分，错误合并和错误拆分都会影响结果\cite{competition}。

\section{方法}
\subsection{全局特征与相似度}
本文分别使用预训练 MegaDescriptor-T-224\cite{mega} 和 MiewID-msv3\cite{miewid}，对输出进行 $\ell_2$ 归一化。MegaDescriptor 使用默认预处理，MiewID 使用 $440\times440$ 像素输入。令两个单位向量为 $z_i^M,z_i^W$，等权融合为
\begin{equation}
 z_i=[\sqrt{0.5}z_i^M,\sqrt{0.5}z_i^W],\qquad
 s_g(i,j)=\frac{z_i^\top z_j+1}{2}.
\end{equation}
所有计算按物种独立进行。本文直接使用比赛图像，不额外训练或分割动物。

\subsection{局部斑纹证据}
蝾螈具有细小、方向多变的斑纹，全局表征容易把同一个体拆开。对每张图像将最长边缩放到不超过 $600$ 像素，最多提取 $1000$ 个 SIFT 关键点。只对融合特征前 $30$ 个近邻候选对进行匹配，按 Lowe 比值 $0.75$ 计数 $m_{ij}$\cite{sift}。局部证据对全局相似度的修正为
\begin{equation}
 s(i,j)=\min\left\{1,\ s_g(i,j)+0.12\,\mathrm{clip}\!\left(\frac{m_{ij}-25}{55},0,1\right)\right\}.
\end{equation}
未列为候选的图像对保持 $s_g$。蝾螈查询之间和查询到参考的匹配都使用修正；角蜥测试图像也计算 SIFT 候选对。猞猁的探索性候选对实验中，SIFT 配对 AUC 低于全局特征，故最终不对猞猁加局部修正。蠵龟使用纯融合特征。

\subsection{测试图片聚类}
对同一物种的测试图片计算距离 $d(i,j)=2[1-s(i,j)]$，使用平均连接层次聚类构造树。通过物种阈值 $\tau_c$ 切树，得到初始测试簇，不预设个体数量。另测试互近邻图方法：每张图像保留 $k=10$ 个最近邻，仅当双向近邻且相似度达到阈值时连边，再取连通分量。图方法作为消融实验，不用于最终 private ARI 最优的提交。

\subsection{已知身份关联}
对测试簇 $C$ 与参考身份 $k$，定义
\begin{equation}
 a(C,k)=\underset{i\in C}{\operatorname{median}}
 \left(\max_{r\in R_k}s(i,r)\right).
\end{equation}
若最高分达到绝对阈值 $\tau_a$，且比第二高分至少高 $0.03$，则将测试簇关联到该已知身份；否则视作新个体簇。德州角蜥没有参考身份，因此跳过此步骤。这里的关联只用于让同一已知个体的多个测试簇合并，提交文件仍只包含测试图片的簇标签。

\section{实验设计}
\subsection{数据与验证协议}
数据来自 AnimalCLEF 2026 比赛页面\cite{competition}，包含 13074 张带标签训练图及 2409 张测试图。猞猁、蝾螈、蠵龟分别有 $2957/946$、$1388/689$、$8729/500$ 张训练/测试图；德州角蜥只有 $274$ 张测试图。每个带标签物种的身份随机分为 75\% 参数校准池与 25\% 保留池。各池中约 60\% 至少有两张图片的身份被模拟为已知个体：一半图片作为参考，另一半作为查询。其他身份只作为查询，模拟新个体。只在校准池搜索 $\tau_c\in\{0.55,0.575,\ldots,0.95\}$ 与 $\tau_a\in\{0.65,0.70,\ldots,0.95\}$，在身份不重合的保留池计算 ARI。使用 2026、2027、2028 三个随机种子。

局部匹配探索性实验从 200 张蝾螈训练图像出发，各取融合特征前 10 个近邻，形成 1941 个不同候选对，其中 91 对同身份。这是偏向相似图像的候选分布，仅用于判断 SIFT 是否值得完整验证。德州角蜥缺乏带标签样本，初始聚类阈值取其他三个物种阈值的中位数 $0.725$，只产生 6 个大簇；SIFT 候选对中高匹配数稀疏，因此比较更保守的 $0.825$、$0.85$、$0.875$ 三个阈值。另报告前三个物种保留查询图片合并后的 ARI，但它不包含角蜥，不能视为官方成绩。

\subsection{迭代准则}
分别评估特征模型、等权融合、局部匹配和图聚类，每次记录三种划分的保留集 ARI、错误合并及错误拆分。角蜥没有独立标签，阈值比较借助 Kaggle 分数，因此单独标出其排行榜适配风险。本文未运行动物分割或微调实验，也不将这些方法计入实验结果。

\section{结果与讨论}
\begin{table}[ht]
\centering
\caption{三个身份互斥保留集的 ARI，均值 $\pm$ 样本标准差。末列不含角蜥。}
\small
\begin{tabular}{lcccc}
\toprule
方法 & 猞猁 & 蝾螈 & 蠵龟 & 合并 \\
\midrule
MegaDescriptor & $.251\pm.165$ & $.097\pm.010$ & $.496\pm.156$ & $.422\pm.139$ \\
MiewID & $.264\pm.092$ & $.152\pm.018$ & $.775\pm.047$ & $.637\pm.058$ \\
等权融合 & $.256\pm.137$ & $.162\pm.024$ & $.835\pm.050$ & $.648\pm.038$ \\
融合 + 互近邻图 & $.278\pm.188$ & $.113\pm.037$ & $.742\pm.047$ & $.619\pm.038$ \\
融合 + 蝾螈 SIFT & $.256\pm.137$ & $.222\pm.023$ & $.835\pm.050$ & $.648\pm.038$ \\
\bottomrule
\end{tabular}
\end{table}

蝾螈候选对的全局特征 AUC 为 $0.735$，SIFT 匹配数 AUC 为 $0.885$，与完整保留集的改进方向一致。融合显著改善蠵龟，但猞猁三个划分的波动很大。图聚类给猞猁带来小幅均值收益，却降低蝾螈、蠵龟的结果。

\begin{table}[ht]
\centering
\caption{Kaggle 提交 ARI。private 为比赛隐藏划分评分，不是本地保留集。}
\small
\begin{tabular}{lccc}
\toprule
提交方法 & 角蜥阈值 & public & private \\
\midrule
MegaDescriptor & $.800$ & $.13528$ & $.18564$ \\
MiewID & $.800$ & $.19256$ & $.25711$ \\
等权融合 & $.725$ & $.21143$ & $.26174$ \\
融合 + 蝾螈 SIFT & $.725$ & $.22721$ & $.27104$ \\
角蜥阈值 $.85$，无角蜥 SIFT & $.850$ & $.26093$ & $.34872$ \\
再加角蜥 SIFT & $.825$ & $.29224$ & $.36086$ \\
再加角蜥 SIFT（最终） & $.850$ & $\mathbf{.30642}$ & $\mathbf{.39072}$ \\
再加角蜥 SIFT & $.875$ & $.29795$ & $.39003$ \\
再加猞猁图聚类 & $.850$ & $.31320$ & $.37797$ \\
\bottomrule
\end{tabular}
\end{table}

角蜥在最终阈值 $.85$ 下被划为 $214$ 个测试簇，原迁移阈值只产生 $6$ 簇。单独提高阈值而不加入角蜥 SIFT，private ARI 已从 $.27104$ 升到 $.34872$；再加入局部匹配，升至 $.39072$。$0.85$ 与 $0.875$ 的 private 分数接近，$0.825$ 则较低。猞猁改用图聚类虽将 public ARI 提到 $.31320$，private ARI 降至 $.37797$，所以最终保留平均连接聚类。这些比较也说明不应把公开榜单的单次提升当作稳健泛化证据。

\section{可复现性与限制}
代码、依赖、随机种子、阈值及实验记录保存在\href{https://github.com/Fangwenky/animalclef-2026-reid}{公开仓库}。比赛图像和标签不上传。验证集通过已标注训练数据构造，其个体比例及拍摄分布未必等同隐藏测试集；角蜥没有直接的本地标签验证，阈值又借助 Kaggle 分数选择，存在排行榜适配。实际 Kaggle 成绩与本地保留集 ARI 分开报告。预训练模型数据与比赛图像是否存在重叠未被独立排查。最终 private ARI $.39072$ 仍明显低于理想划分；后续应优先尝试动物裁剪、局部几何一致性和只用训练数据完成的物种化验证。

\begin{thebibliography}{9}
\bibitem{competition} Adam, L. et al. AnimalCLEF26 @ CVPR \& CLEF. Kaggle, 2026. \href{https://www.kaggle.com/competitions/animal-clef-2026}{比赛页面}.
\bibitem{mega} \v{C}erm\'ak, V. et al. MegaDescriptor-T-224 model card. Hugging Face. \url{https://huggingface.co/BVRA/MegaDescriptor-T-224}.
\bibitem{miewid} Conservation X Labs. MiewID-msv3 model card. Hugging Face. \url{https://huggingface.co/conservationxlabs/miewid-msv3}.
\bibitem{sift} Lowe, D. G. Distinctive image features from scale-invariant keypoints. International Journal of Computer Vision, 60, 91--110, 2004.
\end{thebibliography}
\end{document}
